\section{Ramsey Numbers}

Everything so far has asked when a coloring exists. The opposite question is at least as natural: how large does a structure have to be before every coloring of it is forced to contain something monochromatic? Let's step back into ``graph land'' from ``hypergraph land'' and ask it of the complete graphs, coloring edges this time rather than vertices.

\begin{boxdefinition}[Ramsey Number]
    The $k$th Ramsey number, denoted $R(k)$, is the smallest $n \in \N$ such that any $2$-colouring of edges of $K_n$ gives a monochromatic $K_k$.
\end{boxdefinition}

One immediate result we can see is that $R(3) \geq 6$, which we can see simply by drawing a picture of $K_5$ and studying what happens if we move our attention to $K_6$.

\begin{figure}[H]
    \centering
    \begin{tikzpicture}[scale = 2]
        \begin{scope}[red]
            \stepedges{5}{1}
        \end{scope}
        \begin{scope}[blue]
            \stepedges{5}{2}
        \end{scope}
        \cyclevertices{5}
    \end{tikzpicture}
\end{figure}

Each colour carries a $5$-cycle---the outer pentagon in red, the inner pentagram in blue---and a $5$-cycle contains no triangle at all, so neither colour gives us a monochromatic $K_3$. Nothing forces a monochromatic triangle at $n = 5$, and so $R(3) \geq 6$. The picture also shows how little room there is to spare: every vertex of $K_5$ meets exactly two edges of each colour, and the moment we pass to $K_6$ each vertex has five edges to distribute between two colours, so one of the colours is forced to take three of them.

We can actually prove that $R(3) = 6$, but this is rather nontrivial.

Other known Ramsey numbers are $R(2) = 2$ and $R(4) = 18$. The rest are unknown, though we do have bounds for them. The best (or close to the best) known bounds are
\begin{align*}
    2^{\frac{k}{2}} \leq R(k) \leq 3.8^k
\end{align*}
